\section{Revisão da Literatura}
\label{sec:literatura}

\subsection{Retornos privados e custos da conversão}

A expansão agropecuária envolve uma escolha econômica entre usos alternativos da terra, na qual a conversão da vegetação pode ampliar a renda privada associada à produção, ao mesmo tempo em que reduz estoques de capital natural e serviços ecossistêmicos que nem sempre são incorporados aos retornos percebidos pelos agentes. Seus efeitos podem, portanto, ultrapassar os limites da propriedade convertida. No Cerrado brasileiro, por exemplo, \citeonline{Schwamback2024} documentam, em parcelas experimentais, maior escoamento superficial sob pastagem e maior perda de solo sob cana, em comparação com cerrado arborizado. A comparação refere-se aos usos monitorados no experimento, e não à média do bioma. Na fronteira Amazônia--Cerrado, a substituição de floresta por soja esteve associada à redução do carbono do solo até 100 cm de profundidade, com respostas diferentes conforme a profundidade e o manejo \cite{Popin2025}.

Tais efeitos ajudam a exemplificar por que o resultado produtivo imediato não esgota o valor econômico dos recursos ambientais. \citeonline{Brady2015} propõem valorar serviços do solo a partir de seus efeitos sobre fluxos futuros de renda agrícola, utilizando experimentos de longo prazo na Europa e nos Estados Unidos. O estudo relaciona a depreciação desse capital natural à redução da produtividade e da eficiência do uso de fertilizantes. Sua contribuição é conceitualmente pertinente à tese, embora os parâmetros estimados nesses sistemas não possam ser transferidos diretamente para Mato Grosso.

Com efeito, o problema permanece ainda menos resolvido quando se comparam benefícios e custos privados, como o lucro de um fazendeiro, e públicos, como os impactos gerados para outros agentes e territórios. Isso ocorre porque os serviços ambientais diferem entre lugares, operam em escalas distintas e muitos não possuem preços diretamente observáveis. Estudos que integram informações econômicas e biofísicas mostram a importância da localização e das alternativas de uso da terra para valorar esses serviços e orientar a conservação \cite{Crossman2009,Gascoigne2011}. Assim, a oposição entre conservar e produzir pode refletir diferenças entre retornos apropriados privadamente e custos distribuídos socialmente, além de eventuais conflitos tecnológicos.

Para a Amazônia brasileira, \citeonline{Mann2012} estimam esse descompasso ao comparar o retorno econômico potencial da conversão para a produção de soja com o valor monetário dos serviços ecossistêmicos mantidos pela vegetação nativa. Para tanto, combinam estimativas espacialmente explícitas obtidas a partir de bases, modelos e técnicas de estimação distintos. Sob as hipóteses adotadas, o valor dos serviços supera a renda da soja em parcela expressiva da área analisada, e a omissão desses serviços reduz o bem-estar social relativamente ao ótimo modelado. Esse resultado constitui uma comparação de cenários, não uma estimativa do lucro efetivamente observado em todas as propriedades da região.

No caso do Pantanal, existem tentativas de mensurar economicamente os serviços ambientais, mas os próprios resultados evidenciam dificuldades de escala, cobertura e comparabilidade. \citeonline{SeidlMoraes2000}, ao recalcularem para a sub-região da Nhecolândia estimativas globais com dados locais, obtiveram valor anual por hectare aproximadamente metade daquele derivado dos parâmetros globais, diferença atribuída às condições biofísicas regionais. Mais recentemente, \citeonline{FairbrassEtAl2025}, ao estruturarem indicadores de capital natural para o Pantanal, identificaram lacunas na mensuração da condição e do valor econômico do bioma, das emissões associadas às atividades produtivas e de serviços como ciclagem de nutrientes e biodiversidade.

Além disso, esforços de harmonização entre biomas precisam lidar com componentes de valoração que não são comparáveis. \citeonline{CampoliStivali2026} apresentam uma revisão e harmonização de estimativas para os seis biomas brasileiros, excluindo componentes incompatíveis e produzindo estimativas adicionais para lacunas de informação. Desta feita, a comparação entre biomas exige atenção àquilo que foi efetivamente medido. Este artigo não realizará uma valoração completa dos serviços ambientais; esses estudos justificam investigar como os incentivos privados afetam o estoque que os fornece.

\subsection{Preços, capacidade produtiva e usos alternativos da terra}

Ao largo disso, conforme suscitado logo no início desta seção, a conversão da vegetação nativa pode ser analisada como resultado de escolhas econômicas condicionadas por retornos relativos. \citeonline{Angelsen2010} organiza as políticas de redução do desmatamento em torno das rendas relativas da agricultura e da floresta, distinguindo intervenções que reduzem a atratividade da expansão, elevam os benefícios apropriáveis da conservação ou restringem diretamente a conversão. Esse enquadramento permite compreender por que a intensificação produtiva pode ampliar a pressão local de expansão em determinadas condições, sem presumir que produzir mais implique necessariamente converter mais terra.

A evidência de \citeonline{LubowskiEtAl2008}, para os Estados Unidos, relaciona os retornos líquidos dos usos alternativos e a qualidade da terra às decisões dos proprietários. Na Amazônia brasileira, \citeonline{Pfaff1999} identifica a relevância das características da terra e dos custos de transporte, advertindo que os efeitos estimados de políticas podem refletir a própria seleção de onde elas são implementadas. \citeonline{ChomitzThomas2003}, por sua vez, analisam conversão e intensidade pecuária em escala espacial detalhada, considerando precipitação, solos, infraestrutura, acesso a mercados e proteção territorial. Essas contribuições sustentam a inclusão da capacidade produtiva e das condições locais na análise de preços.

Em escala tropical, \citeonline{BermanEtAl2023} combinam preços internacionais e aptidão agrícola para examinar a relação entre incentivos de mercado e desmatamento. O desenho é especialmente pertinente à H1-a porque um mesmo movimento de preço encontra diferentes possibilidades produtivas no espaço. Os autores também discutem a possibilidade de os preços refletirem choques de oferta e a necessidade de justificar a exogeneidade das exposições. A disponibilidade de uma série internacional, portanto, não elimina por si só os problemas de identificação.

Sem embargo, a relação entre preços e conversão depende do uso que se expande e daquele que é substituído. Na bacia do Tapajós, \citeonline{Braganca2018} encontra que a elevação do preço relativo de lavouras em relação à carne favoreceu a substituição de pastagens por cultivos (\textit{pasture-to-cropland}) e reduziu o desmatamento local. Uma mudança nos preços relativos entre atividades, portanto, não se confunde com uma elevação do retorno esperado da conversão da vegetação nativa. A pressão para abrir novas áreas agropastoris deve, assim, ser distinguida da reorganização da produção em áreas já convertidas. Para este artigo, essas margens de ajuste devem também ser distinguidas da intensificação produtiva. Soja, milho e pecuária tampouco podem ser combinados em um único índice sem explicitar unidades, exposição e mecanismo econômico.

\subsection{Instituições, implementação e respostas aos incentivos}

Assim sendo, essa distinção também importa para a avaliação das políticas de conservação, pois seus incentivos e restrições podem modificar tanto a pressão para abrir novas áreas quanto as escolhas produtivas em áreas já convertidas. A classificação de políticas e mecanismos de conservação como instrumentos regulatórios ou econômicos, por si só, não esclarece como os produtores respondem a essas intervenções. Deste modo, a resposta dos produtores depende de como os incentivos e as restrições modificam os retornos dos usos alternativos da terra, das possibilidades de ajuste da produção e das condições locais de implementação \cite{Angelsen2010}.

Sob essa perspectiva, o enquadramento econômico não supõe efeitos uniformes das instituições. \citeonline{PachecoMeyer2022} comparam regimes fundiários em diferentes contextos brasileiros e mostram que os resultados dependem do regime de origem, da alternativa de destinação da terra e do contexto socioambiental. Já \citeonline{ProbstEtAl2020}, ao examinar o programa Terra Legal, encontram aumento do desmatamento entre pequenos e médios proprietários titulados e maior resposta aos preços agropecuários. A regularização, desse modo, não deve ser confundida com um indicador inequívoco de conservação.

A mesma cautela se aplica ao Cadastro Ambiental Rural. \citeonline{JungEtAl2022} analisam Pará e Mato Grosso e identificam respostas distintas da conversão para lavouras e pastagens, com heterogeneidade por porte dos estabelecimentos e conformidade ambiental. Os resultados apontam o acesso ao crédito como um possível mecanismo. Essa evidência reforça a necessidade de examinar como instrumentos distintos afetam oportunidades produtivas e custos de conversão, em vez de reunir todo registro administrativo sob uma medida única de qualidade institucional.

Fora do contexto brasileiro, \citeonline{ViallonEtAl2019} mostram, em doze casos de uso da terra na Suíça, que a implementação local depende da capacidade das autoridades e dos direitos adquiridos pelos proprietários. \citeonline{AnderssonEtAl2014}, a partir de comunidades florestais na Bolívia, relacionam condições das florestas a diferentes combinações de elaboração de regras, monitoramento e sanções organizadas pelos usuários. Esses trabalhos ajudam a motivar a dimensão territorial da análise; não identificam, para Mato Grosso, a superioridade causal de um nível de governança sobre outro.

Isso não exclui, por óbvio, a relevância de instrumentos econômicos de conservação, como pagamentos por serviços ambientais, mercados de carbono ou financiamento ambiental. As receitas potenciais associadas à conservação do carbono florestal, por exemplo, podem modificar os retornos relativos enfrentados pelos agentes \cite{Angelsen2010,SticklerEtAl2013}. O recorte deste artigo, contudo, concentra-se nos preços agropecuários e nas condições institucionais que poderão ser observadas em séries territoriais comparáveis. Nenhuma conclusão sobre a eficácia de um instrumento específico será extraída sem observar sua aplicação e dispor de um desenho de avaliação apropriado.

\subsection{Fiscalização e pressão econômica em Mato Grosso}

A evidência disponível reconhece que monitoramento, sanções críveis e restrições legais podem conter a conversão, embora seus resultados variem conforme a pressão econômica. \citeonline{HardingEtAl2021} investigam três instrumentos na Amazônia brasileira, entre os quais a Moratória da Soja, que atua por critérios de compra estabelecidos pelas empresas da cadeia produtiva. Encontram redução do desmatamento associado aos preços nos municípios da lista prioritária, maior sensibilidade a preços de outros cultivos em municípios expostos à Moratória da Soja e ampliação da pressão sobre terras não protegidas associada às zonas de conservação. Os resultados sugerem observar também substituições e deslocamentos espaciais, em vez de limitar a avaliação ao efeito médio na área diretamente alcançada pela política.

\citeonline{AssuncaoEtAl2023}, por sua vez, utilizam limitações de detecção por nuvens no sistema DETER como fonte de variação para identificar a atuação fiscalizatória. O estudo encontra efeito de contenção do desmatamento pelo monitoramento e pela fiscalização. Seu desenho evidencia por que a correlação entre número de autos e conversão é insuficiente: a fiscalização pode se concentrar justamente onde a pressão de desmatamento é maior. A estratégia com nuvens tampouco será importada automaticamente para este artigo, pois condições meteorológicas podem afetar a agricultura e a conversão por outras vias.

Em Mato Grosso, \citeonline{Sousa2016} examina multas, embargos e apreensões por meio de regressões e análise espacial, reportando correlações negativas com o desmatamento e diferenças de implementação entre sanções. O estudo destaca a baixa arrecadação das multas e a relevância dos embargos. Esses resultados constituem evidência sobre associações e funcionamento administrativo, sem dispensar o enfrentamento da seleção da fiscalização. Por sua vez, \citeonline{SticklerEtAl2013} analisam custos e consequências potenciais do cumprimento das obrigações de conservação previstas no Código Florestal no estado e discutem a instabilidade regulatória e a combinação com incentivos positivos à conservação.

Esses resultados não são propriamente incompatíveis: refletem desenhos voltados a diferentes margens da decisão de uso da terra, em períodos e ambientes institucionais distintos. A questão que permanece para este artigo é como observar, sob uma base territorial comum, os incentivos de conversão e as condições institucionais que podem ampliar ou restringir a resposta dos agentes. Trata-se de delimitar uma contribuição empírica para Mato Grosso, sem afirmar ausência de estudos que já combinam preços e instituições.

\subsection{Heterogeneidade territorial e contribuição do artigo}

Em Mato Grosso, a expansão de duas safras anuais está associada ao acesso a transportes, à produção agrícola anterior e à conexão com mercados, como mostram \citeonline{VanWeyEtAl2013}. Os mapas de \citeonline{SimoesEtAl2020} permitem distinguir trajetórias de agricultura, pastagem e vegetação entre 2001 e 2017. No Pantanal, \citeonline{SeidlEtAl2001} relacionam a conversão para pastagens cultivadas a indicadores de riqueza fundiária e pecuária, esforço agrícola e localização. Esse estudo abrange o Pantanal brasileiro e não deve ser tratado como estimativa específica dos municípios pantaneiros de Mato Grosso.

Tomadas em conjunto, essas contribuições justificam investigar a sensibilidade da conversão aos preços e seu condicionamento institucional com atenção às atividades, aos estoques remanescentes e aos biomas. A dimensão intertemporal aparece no horizonte em que produtores percebem preços, incorrem em custos de conversão e antecipam a aplicação das regras. A análise proposta observará defasagens e persistência das respostas, mas não estimará preferências temporais ou expectativas individuais a partir dessas defasagens.

Dito isso, a análise que se segue parte da caracterização dessas diferenças no território para, posteriormente, estimar as relações previstas em H1, H1-a e H1-b. Busca-se avançar na compreensão de como condições econômicas e institucionais alteram a formação do estoque ambiental, preservando a distinção entre mecanismos conceituais, associações observadas e efeitos cuja interpretação causal possa ser sustentada.

